\subsection{Задание 2. Взлом шифра Хилла}

\subsubsection{Постановка задачи}
Смоделируем ситуацию, когда у нас есть два зашифрованных сообщения $C_1$ и $C_2$, полученных шифром Хилла с одним и тем же неизвестным ключом $K$ размера $2\times 2$, и известен открытый текст $P_1$, соответствующий $C_1$. Требуется найти $P_2$ — открытый текст второго сообщения.

Используется тот же алфавит из $n = 33$ символов (см. таблицу~\ref{tab:code_tab}), что и в Задании 1.

\subsubsection{Генерация ключа и шифрование}

Случайным образом сгенерируем ключ $K$ размера $2\times 2$ с обратимым по модулю 33 определителем:
\[
K = \begin{pmatrix} 5 & 3 \\ 2 & 7 \end{pmatrix}, \qquad \det(K) = 5\cdot 7 - 3\cdot 2 = 35 - 6 = 29.
\]
Проверка: $\text{НОД}(29, 33) = 1$ — ключ обратим по модулю 33.

\paragraph{Первое сообщение $P_1$ (12 символов):}
\texttt{КРИПТОГРАФИЯ}

Числовые коды (по таблице~\ref{tab:code_tab}):
\[
P_1 = (10,\; 16,\; 8,\; 15,\; 18,\; 14,\; 3,\; 16,\; 0,\; 19,\; 8,\; 30).
\]

\paragraph{Второе сообщение $P_2$ (12 символов):}
\texttt{СЕКРЕТНОЕСОО} (С,Е,К,Р,Е,Т,Н,О,Е,С,О,О).

Числовые коды:
\[
P_2 = (17,\; 5,\; 10,\; 16,\; 5,\; 18,\; 13,\; 14,\; 5,\; 17,\; 14,\; 14).
\]

\paragraph{Шифрование.} Каждые два символа образуют вектор-строку, умножаемую на $K$:
$C_i = (P_i \cdot K) \bmod 33$.

\noindent\textbf{Сообщение $P_1$:}
\[
\begin{aligned}
(10,16) \cdot K &= (10\cdot 5 + 16\cdot 2,\; 10\cdot 3 + 16\cdot 7) = (50+32,\; 30+112) = (82, 142) \equiv (16, 10) \pmod{33},\\
(8,15) \cdot K &= (8\cdot 5 + 15\cdot 2,\; 8\cdot 3 + 15\cdot 7) = (40+30,\; 24+105) = (70, 129) \equiv (4, 30) \pmod{33},\\
(18,14) \cdot K &= (18\cdot 5 + 14\cdot 2,\; 18\cdot 3 + 14\cdot 7) = (90+28,\; 54+98) = (118, 152) \equiv (19, 20) \pmod{33},\\
(3,16) \cdot K &= (3\cdot 5 + 16\cdot 2,\; 3\cdot 3 + 16\cdot 7) = (15+32,\; 9+112) = (47, 121) \equiv (14, 22) \pmod{33},\\
(0,19) \cdot K &= (0\cdot 5 + 19\cdot 2,\; 0\cdot 3 + 19\cdot 7) = (38, 133) \equiv (5, 1) \pmod{33},\\
(8,30) \cdot K &= (8\cdot 5 + 30\cdot 2,\; 8\cdot 3 + 30\cdot 7) = (40+60,\; 24+210) = (100, 234) \equiv (1, 3) \pmod{33}.
\end{aligned}
\]
Шифртекст $C_1$: коды $(16,10,4,30,19,20,14,22,5,1,1,3)$, что соответствует строке
\[
C_1 = \texttt{РКАЬУФОХЖББГ}.
\]

\noindent\textbf{Сообщение $P_2$:}
\[
\begin{aligned}
(17,5) \cdot K &= (17\cdot 5 + 5\cdot 2,\; 17\cdot 3 + 5\cdot 7) = (85+10,\; 51+35) = (95, 86) \equiv (29, 20) \pmod{33},\\
(10,16) \cdot K &= (82, 142) \equiv (16, 10) \pmod{33},\\
(5,18) \cdot K &= (5\cdot 5 + 18\cdot 2,\; 5\cdot 3 + 18\cdot 7) = (25+36,\; 15+126) = (61, 141) \equiv (28, 9) \pmod{33},\\
(13,14) \cdot K &= (13\cdot 5 + 14\cdot 2,\; 13\cdot 3 + 14\cdot 7) = (65+28,\; 39+98) = (93, 137) \equiv (27, 5) \pmod{33},\\
(5,17) \cdot K &= (5\cdot 5 + 17\cdot 2,\; 5\cdot 3 + 17\cdot 7) = (25+34,\; 15+119) = (59, 134) \equiv (26, 2) \pmod{33},\\
(14,14) \cdot K &= (14\cdot 5 + 14\cdot 2,\; 14\cdot 3 + 14\cdot 7) = (70+28,\; 42+98) = (98, 140) \equiv (32, 8) \pmod{33}.
\end{aligned}
\]
Шифртекст $C_2$: коды $(29,20,16,10,28,9,27,5,26,2,32,8)$, что соответствует
\[
C_2 = \texttt{ЭУРКЬЙЫЖЪВЯИ}.
\]

\subsubsection{Атака: восстановление ключа по известной паре $(P_1, C_1)$}

Пусть ключ имеет вид $K = \begin{pmatrix} a & b \\ c & d \end{pmatrix}$. Уравнение $P_1 \cdot K = C_1 \pmod{33}$ даёт 12 уравнений на 4 неизвестных. Возьмём первые два блока (первые 4 символа $P_1$ и $C_1$):

\[
\begin{cases}
10 a + 16 c \equiv 16 \pmod{33} \\
10 b + 16 d \equiv 10 \pmod{33} \\
8 a + 15 c \equiv 4 \pmod{33} \\
8 b + 15 d \equiv 30 \pmod{33}
\end{cases}
\]

Из первой пары (для $a, c$):
\[
\begin{cases}
10 a + 16 c \equiv 16 \\
8 a + 15 c \equiv 4
\end{cases}
\]

Решим через определитель:

Матрица системы для $(a, c)$:
\[
M = \begin{pmatrix} 10 & 16 \\ 8 & 15 \end{pmatrix}, \qquad \det(M) = 10\cdot 15 - 16\cdot 8 = 150 - 128 = 22.
\]
Проверка: $\text{НОД}(22, 33) = 11 \neq 1$ — матрица необратима! Значит, этот блок не подходит для прямого обращения.

Возьмём другой блок. Попробуем блоки 3–4 (символы 5–8): $P_1$: $(18,14)$ и $(3,16)$; $C_1$: $(19,20)$ и $(14,22)$.

\[
\begin{cases}
18 a + 14 c \equiv 19 \\
3 a + 16 c \equiv 14
\end{cases}
\]
\[
M = \begin{pmatrix} 18 & 14 \\ 3 & 16 \end{pmatrix}, \qquad \det(M) = 18\cdot 16 - 14\cdot 3 = 288 - 42 = 246 \equiv 246 - 7\cdot 33 = 246 - 231 = 15 \pmod{33}.
\]
$\text{НОД}(15, 33) = 3 \neq 1$ — снова необратима.

Возьмём блоки 1 и 5 (символы 1–2 и 9–10): $P_1$: $(10,16)$ и $(0,19)$; $C_1$: $(16,10)$ и $(5,1)$.

\[
\begin{cases}
10 a + 16 c \equiv 16 \\
0\cdot a + 19 c \equiv 5
\end{cases}
\]
Из второго: $19 c \equiv 5 \pmod{33}$. Найдём $19^{-1} \bmod 33$. $\text{НОД}(19,33)=1$.
$33 = 1\cdot 19 + 14$, $19 = 1\cdot 14 + 5$, $14 = 2\cdot 5 + 4$, $5 = 1\cdot 4 + 1$.
$1 = 5 - 4 = 5 - (14 - 2\cdot 5) = 3\cdot 5 - 14 = 3\cdot(19-14) - 14 = 3\cdot 19 - 4\cdot 14 = 3\cdot 19 - 4\cdot(33 - 19) = 7\cdot 19 - 4\cdot 33$.
Значит $19^{-1} \equiv 7 \pmod{33}$.
$c \equiv 7\cdot 5 = 35 \equiv 2 \pmod{33}$.

Подставим в первое: $10 a + 16\cdot 2 \equiv 16 \Rightarrow 10 a + 32 \equiv 16 \Rightarrow 10 a \equiv -16 \equiv 17 \pmod{33}$.
$10^{-1} \bmod 33$: $\text{НОД}(10,33)=1$. $33 = 3\cdot 10 + 3$, $10 = 3\cdot 3 + 1$, $1 = 10 - 3\cdot 3 = 10 - 3(33 - 3\cdot 10) = 10\cdot 10 - 3\cdot 33$. Значит $10^{-1} \equiv 10$.
$a \equiv 10\cdot 17 = 170 \equiv 170 - 5\cdot 33 = 170 - 165 = 5 \pmod{33}$.

Итак $a = 5$, $c = 2$. Аналогично для $(b,d)$ из уравнений:
\[
\begin{cases}
10 b + 16 d \equiv 10 \\
19 d \equiv 1
\end{cases}
\]
$d \equiv 7\cdot 1 = 7 \pmod{33}$.
$10 b + 16\cdot 7 \equiv 10 \Rightarrow 10 b + 112 \equiv 10 \Rightarrow 10 b \equiv 10 - 112 = -102 \equiv -102 + 4\cdot 33 = -102 + 132 = 30 \pmod{33}$.
$b \equiv 10\cdot 30 = 300 \equiv 300 - 9\cdot 33 = 300 - 297 = 3 \pmod{33}$.

\textbf{Найденный ключ:}
\[
K = \begin{pmatrix} 5 & 3 \\ 2 & 7 \end{pmatrix},
\]
что полностью совпадает с исходным ключом. Проверим на всех блоках $P_1$:

$(10,16)K \equiv (16,10)$, $(8,15)K \equiv (4,30)$, $(18,14)K \equiv (19,20)$ , $(3,16)K \equiv (14,22)$ , $(0,19)K \equiv (5,1)$ , $(8,30)K \equiv (1,3)$ .

\subsubsection{Нахождение обратной матрицы $K^{-1}$}

$\det(K) = 29$. Найдём $29^{-1} \bmod 33$.
$33 = 1\cdot 29 + 4$, $29 = 7\cdot 4 + 1$, $1 = 29 - 7\cdot 4 = 29 - 7(33 - 29) = 8\cdot 29 - 7\cdot 33$. Значит $29^{-1} \equiv 8 \pmod{33}$.

Алгебраические дополнения: $A_{11} = 7$, $A_{12} = -2$, $A_{21} = -3$, $A_{22} = 5$.
\[
A^T = \begin{pmatrix} 7 & -2 \\ -3 & 5 \end{pmatrix} \equiv \begin{pmatrix} 7 & 31 \\ 30 & 5 \end{pmatrix} \pmod{33}.
\]
\[
K^{-1} = 8 \cdot \begin{pmatrix} 7 & 31 \\ 30 & 5 \end{pmatrix} = \begin{pmatrix} 56 & 248 \\ 240 & 40 \end{pmatrix} \equiv \begin{pmatrix} 23 & 17 \\ 9 & 7 \end{pmatrix} \pmod{33}.
\]

Проверка:
\[
K \cdot K^{-1} = \begin{pmatrix} 5 & 3 \\ 2 & 7 \end{pmatrix} \begin{pmatrix} 23 & 9 \\ 17 & 7 \end{pmatrix} = \begin{pmatrix} 115+51 & 45+21 \\ 46+119 & 18+49 \end{pmatrix} = \begin{pmatrix} 166 & 66 \\ 165 & 67 \end{pmatrix} \equiv \begin{pmatrix} 1 & 0 \\ 0 & 1 \end{pmatrix} \pmod{33}.
\]

\subsubsection{Расшифровка второго сообщения $C_2$}

$C_2 = (29,20,16,10,28,9,27,5,26,2,32,8)$. Вычисляем $P_2 = C_2 \cdot K^{-1} \bmod 33$:

\[
\begin{aligned}
(29,20) \cdot K^{-1} &= (29\cdot 23 + 20\cdot 17,\; 29\cdot 9 + 20\cdot 7) = (667+340,\; 261+140) = (1007, 401) \\
&\equiv (1007 - 30\cdot 33,\; 401 - 12\cdot 33) = (1007-990,\; 401-396) = (17, 5) \pmod{33},\\
(16,10) \cdot K^{-1} &= (16\cdot 23 + 10\cdot 17,\; 16\cdot 9 + 10\cdot 7) = (368+170,\; 144+70) = (538, 214) \\
&\equiv (538 - 16\cdot 33,\; 214 - 6\cdot 33) = (538-528,\; 214-198) = (10, 16) \pmod{33},\\
(28,9) \cdot K^{-1} &= (28\cdot 23 + 9\cdot 17,\; 28\cdot 9 + 9\cdot 7) = (644+153,\; 252+63) = (797, 315) \\
&\equiv (797 - 24\cdot 33,\; 315 - 9\cdot 33) = (797-792,\; 315-297) = (5, 18) \pmod{33},\\
(27,5) \cdot K^{-1} &= (27\cdot 23 + 5\cdot 17,\; 27\cdot 9 + 5\cdot 7) = (621+85,\; 243+35) = (706, 278) \\
&\equiv (706 - 21\cdot 33,\; 278 - 8\cdot 33) = (706-693,\; 278-264) = (13, 14) \pmod{33},\\
(26,2) \cdot K^{-1} &= (26\cdot 23 + 2\cdot 17,\; 26\cdot 9 + 2\cdot 7) = (598+34,\; 234+14) = (632, 248) \\
&\equiv (632 - 19\cdot 33,\; 248 - 7\cdot 33) = (632-627,\; 248-231) = (5, 17) \pmod{33},\\
(32,8) \cdot K^{-1} &= (32\cdot 23 + 8\cdot 17,\; 32\cdot 9 + 8\cdot 7) = (736+136,\; 288+56) = (872, 344) \\
&\equiv (872 - 26\cdot 33,\; 344 - 10\cdot 33) = (872-858,\; 344-330) = (14, 14) \pmod{33}.
\end{aligned}
\]

Полученные коды: $(17,5,10,16,5,18,13,14,5,17,14,14)$, что соответствует слову
\[
P_2 = \texttt{СЕКРЕТНОЕСОО}.
\]

\textbf{Проверка первого сообщения:} $P_1 = C_1 \cdot K^{-1}$:
\[
(16,10)K^{-1} \equiv (10,16),\; (4,30)K^{-1} \equiv (8,15),\; (19,20)K^{-1} \equiv (18,14),\;
(14,22)K^{-1} \equiv (3,16),\;
(5,1)K^{-1} \equiv (0,19),\;
(1,3)K^{-1} \equiv (8,30).
\]
Восстановлено $P_1 = \texttt{КРИПТОГРАФИЯ}$.

\subsubsection{Вывод по Заданию 2}

Атака на основе известного открытого текста на шифр Хилла оказывается успешной, потому что шифр линеен. Зная одну пару (открытый текст, шифртекст), можно составить систему линейных уравнений относительно элементов ключевой матрицы и решить её по модулю $n$. Если какой-то блок матрицы $P_1$ оказывается необратимым по модулю 33, это не мешает атаке — достаточно взять другой блок или решить систему целиком.

\textbf{Основные выводы:}
\begin{enumerate}
    \item Шифр Хилла уязвим к атаке на основе известного открытого текста: одного сообщения из 12 символов достаточно для полного восстановления ключа $2\times 2$ (нужно всего 4 независимых уравнения).
    \item Для ключей большего размера потребуется больше символов известного открытого текста (для $m\times m$ — порядка $m^2$ символов).
    \item Необратимость отдельных блоков открытого текста по модулю 33 не является препятствием: можно использовать другие блоки или решать систему совместно.
    \item Увеличение размера ключа повышает стойкость, но принципиально не защищает от атаки на основе известного открытого текста.
\end{enumerate}